
\documentclass[11pt]{article}

% Anonymous review version for ARR / EMNLP 2026.
\usepackage[review]{acl}

\usepackage{times}
\usepackage{latexsym}
\usepackage[T1]{fontenc}
\usepackage[utf8]{inputenc}
\usepackage{microtype}
\usepackage{inconsolata, multirow}
\usepackage{booktabs}
\usepackage{array}
\usepackage{graphicx}
% \usepackage{paralist}
\usepackage{enumitem}
\usepackage{listings}
\input{math_commands.tex}
\usepackage{cleveref}
\usepackage{placeins}
\usepackage{tabularx}

\makeatletter
\setlength{\@dblfptop}{0pt}
\makeatother

\lstdefinestyle{prompt}{
  language=,                % treat as plain text
  basicstyle=\ttfamily\scriptsize,
  breaklines=true,          % wrap long lines
  breakindent=0pt,
  columns=fullflexible,
  backgroundcolor=\color{black!3},
  frame=single,
  framerule=0.4pt,
  rulecolor=\color{black!40},
  xleftmargin=0pt,
  xrightmargin=0pt,
  aboveskip=4pt,
  belowskip=4pt,
}

\graphicspath{{figures/}}

\definecolor{darkblue}{rgb}{0, 0, 0.5}
\hypersetup{colorlinks=true, citecolor=darkblue, linkcolor=darkblue, urlcolor=darkblue}

\title{Knowledge Acquisition During Pre-training?\\ Large Language Models Learn Better With Auxiliary Views}

\author{Anonymous ARR submission}

\newcommand{\fix}{\marginpar{FIX}}
\newcommand{\new}{\marginpar{NEW}}

\begin{document}

\maketitle

\begin{abstract}
Gaps remain in our understanding of how large language models (LLMs) acquire knowledge during pre-training. We posit that auxiliary views, reformulations of knowledge across different contexts, are significantly helpful for learning. We design controlled experiments to isolate this effect. First, we confirm that repetition is necessary for acquisition and clarify that paraphrasing helps, but only at smaller batch sizes. Second, holding the token budget fixed, we find that reallocating tokens from document repetition to auxiliary views improves knowledge acquisition, even factual recall. Third, we identify forms of knowledge, contextual and foundational, that aid learning in the presence of prior knowledge gaps. Finally, we examine how these effects arise mechanistically, via layer-wise biases in encoding. Together, our findings suggest that auxiliary representations of knowledge, which arise naturally in large pre-training corpora, are a key factor in the success of pre-training, and offer a plausible explanation for why data diversity matters.
\end{abstract}

\section{Introduction}


Various aspects of data have been shown to be important in pre-training, including deduplication~\citep{raffel2020exploring, lee2021deduplicating, zhang2022opt}, filtering~\citep{weber2024redpajama, li2024datacomp}, coverage and depth ~\citep{kandpal2023large}, quality~\citep{gunasekar2023textbooks, longpre2024pretrainer}, and diversity~\citep{chen2025revisiting, zhang2025harnessing}. However, these insights entail corpus characteristics, overlooking a fundamental question: \textit{how should knowledge be represented?}

The complexity of this question can vary. For atomic facts like biographical attributes, representation is relatively trivial. But for complex domains such as biomedicine or law, where knowledge is multifaceted and interdependent, formulation is far less obvious. This is becoming relevant as researchers adapt LLMs to specialized domains via continued pre-training~\citep{bai2025intern,sellergren2025medgemma,wang2025txgemma,luo2025large,colombo2024saullm, singhal2025toward}.
\begin{figure}[t]
    \centering
    \includegraphics[width=\columnwidth]{EMNLP/figures/probe_32b_conditions_2x2_lines.pdf}
\caption{\small{Three training mixes on OLMo2-32B: \textit{Source}, \textit{Para.\ 9}, and \textit{Para.\ 9 + Aux}. Across all metrics, adding paraphrases improves over source-only, and adding auxiliary views improves rather considerably. See Section \ref{aux} for full results.}}
  \label{fig1:probe_32b}
\end{figure}

Thus, for our experiments, we collect recent arXiv papers, legal opinions, and medical case studies to serve as self-contained domain knowledge. We continue pre-training on these texts and related formulations, to study the acquisition of \textit{complex knowledge that builds upon the LLM's existing knowledge}. In this setting, we study two research questions:

\textbf{RQ 1:} How should knowledge be represented in natural language?

\textbf{RQ 2:} What surrounding knowledge should be represented alongside it?


Our experiments reveal three main findings. (1) Acquisition of new knowledge requires repetition and paraphrasing prevents the overfitting that can occur; however, the benefits of paraphrasing diminish with larger batch sizes. (2) Holding the token budget fixed, reallocating tokens from document repetition to auxiliary views improve the learning in both understanding and factual recall, and leaves a distinct signature in how the knowledge is encoded mechanistically. (3) Bridging knowledge gaps by training the LLM on prerequisite and contextual knowledge improves learning. These results lead us to a central conjecture: during pre-training, LLMs benefit from \textit{auxiliary views} (a web of explanations, analogies, and reformulations that humans generate as they learn, process, and teach each other) and acquire a more generalizable encoding of knowledge.


We hope these findings  provide further insight into how language models acquire knowledge, principles to improve continued pre-training recipes on low-data domains, and a deeper explanation to the importance of data diversity.


\section{Related Works}
\label{related_works}


\textbf{Pre-training.} Model size and data volume are the primary determinants of LLMs' capabilities ~\citep{brown2020language, kaplan2020scaling, carlini2022quantifying, tirumala2022memorization}. While model size is usually emphasized, \citet{hoffmann2022training} shows optimal performance requires scaling data together. This large-scaled nature of pre-training has made understanding the role of data challenging. Existing studies, varying in corpora, model architectures and sizes, and evaluation methods, often yield conflicting findings. Even the subject of data repetition has been well-debated~\citep{lee2021deduplicating, taylor2022galactica, hernandez2022scaling, xue2023repeat, muennighoff2023scaling}.  


\textbf{Knowledge Acquisition during Pre-training.} Thus, we build upon previous efforts that study how pre-training instills knowledge in a controlled manner. \citet{allen2024physics} show that paraphrased augmentation increases the memorization of biographical facts from 9.7\% to 96.6\% . They find that incorporating QAs about facts during pre-training leads to increased encoding on held-out facts, whereas inserting QAs during instruction-tuning does not. This strongly suggests the importance of  data formulation in pre-training.

Further, \citet{chang2024large} intermittently injects fictional facts during pre-training, observing that knowledge is acquired incrementally upon each exposure and subject to decay, suggesting a gradual, not emergent, learning process. Confounding factors are controlled to clearly isolate the effect.

A limitation of these prior works is that \citet{chang2024large} and \citet{allen2023physics} only use biographical facts; they also provide conflicting views on the benefits of paraphrasing. We focus on complex knowledge, and offer an explanation for these reported differences.

\textbf{Domain Adaptation via Continued pre-training.} 
Teaching large language models (LLMs) a specific set of new knowledge via continued pre-training is difficult
\citep{wang2021can, jang2021towards, hu2023meta,ovadia2024fine, hoffbauer-etal-2024-knowledge}. For instance, a 70B model continually pretrained on Wiki-style documents, despite sophisticated augmentation, only recalls 62.7\% of facts~\citep{jiang2024instruction}. In practice, continued pre-training is adopted in varying levels. For example, MedGemma~\citep{sellergren2025medgemma} and TxGemma~\citep{wang2025txgemma} do not perform any text-based continued pre-training, while Intern-S1~\citep{bai2025intern} does for 5 trillion tokens. We believe our study provides practical takeaways that can make later-stage pre-training more reliable.

% While other forms of knowledge injection are important ~\citep{ovadia2024fine, song2025injecting}, the scope of our study is pre-training.

\begin{table}[t]
\centering
{\footnotesize
\begin{tabularx}{\linewidth}{@{} >{\raggedright\arraybackslash}X @{}}
\toprule
\textbf{Blog.} Direct Preference Optimization (DPO) offers a neat change of perspective. Instead of treating reward modeling and policy optimization as separate steps, DPO re-parameterizes the reward class so [...] has a closed form. That lets you fit the policy directly to human comparisons with a simple classification loss --- no RL required. \tabularnewline
\midrule
\textbf{Stack Exchange.} \textbf{Q:} How exactly does the partition function cancel out in the DPO derivation? They say the partition function term $\beta \log Z(x)$ cancels --- can someone show the algebraic steps explicitly? \textbf{A:} Short answer: because the partition function term in the reparameterization depends only on the prompt $x$, it is identical for both completions and therefore subtracts out when you form the [...]\tabularnewline
\midrule
\textbf{Textbook.} The partition function term $\beta\log Z(x)$ depends only on $x$, so it cancels when forming differences. Preference models like Bradley--Terry depend only on reward differences, not on absolute reward values. Therefore the unknown normalization that made reward recovery hard is irrelevant to the likelihood of observed pairwise comparisons. This is precisely why we can reparameterize rewards in terms of [...] \tabularnewline
\bottomrule
\end{tabularx}
}
\caption{\small{We rewrite each document into three genres while preserving the underlying knowledge. Excerpts provided above. The \textit{blog} adopts an accessible, narrative tone; the \textit{Stack Exchange} entry entails question-and-answer exchange; and the \textit{textbook} situates it in a formal exposition.}}
\label{tab1}
\vspace{-5mm}
\end{table}

\begin{table*}[t]
\centering
{\footnotesize
\resizebox{0.98\textwidth}{!}{
\begin{tabular}{@{} >{\raggedright\arraybackslash}p{3cm} >{\raggedright\arraybackslash}p{12cm} @{}}
\toprule
\textbf{Original sentence} & The \textit{added constraint is important}, as it \textit{prevents the model from} deviating too far from the distribution [...], as well as \textit{maintaining the generation diversity} and [...] \\
\midrule
\textbf{Factual probe} & In the paper `Direct Preference Optimization [...]', the authors state that the \textit{added constraint} in the reinforcement learning objective not only \textit{prevents the model from} deviating too far from the distribution on which the reward model is accurate, but also \textit{maintains} \textbf{the generation diversity}. \\
\midrule
\textbf{Support sentences} & [...] still \textit{expensive to estimate the partition function} $Z(x)$ [...] the Bradley-Terry model \textit{depends only on the difference of rewards} [...] Substituting the reparameterization [...] into the preference model, \textit{the partition function cancels}, and we can express the human preference probability in terms of only the optimal policy $\pi^*$ and reference policy [...] \\
\midrule
\textbf{Inference probe} & According to the paper `Direct Preference Optimization [...]', the expensive quantity from the optimal-policy form that becomes unnecessary to estimate once DPO rewrites the Bradley-Terry preference model in terms of policies is \textbf{the partition function}. \\
\bottomrule
\end{tabular}
}}
\caption{\small{The \textit{factual} probe extracts questions from knowledge-bearing sentences, converts them into statements with answers at the end, and adds context. The \textit{inference} probe composes knowledge from one or multiple sentences (see support sentences) to infer information that was not explicitly stated before. In this case, it integrates how the $Z(x)$ is expensive to estimate, the Bradley-Terry model depends only on reward differences, and that substituting the reparameterization cancels out to infer that the partition function is the quantity that no longer needs to be estimated. The \textbf{target span} of each probe is bolded.}}
\label{tab2}
\vspace{-5mm}
\end{table*}

\section{Experimental Setup}
\label{experiment}
 
\textbf{Problem Formulation.}
The pre-training objective for an autoregressive LM, $f$ parameterized by $\theta$, is next-token prediction. 
Given a corpus $C$ of documents ${d_1,\dots,d_M}$, each a sequence of tokens $(t_{m,1}, \dots, t_{m,n_m})$,
we minimize the loss function:

\begin{equation}
L(\theta) = -\sum_{m=1}^{M} \sum_{i=1}^{n_m} \log P(t_{m,i} \mid t_{m,<i}; \theta)
\end{equation}

We denote knowledge abstractly as $K$, accessible through the sub-corpus $\mathcal{C}_K = \{\, d_i \in \mathcal{C} \mid d_i \text{ contains information about } K \,\}$. Because $K$ is acquired only through next-token prediction over $\mathcal{C}_K$, the selection of tokens should matter. This motivates our central question: how should knowledge be represented in text, and how should we construct the documents $d \in \mathcal{C}_K$ to better instill $K$?


\textbf{Observation}. Knowledge is rarely represented once in a corpus. Consider ``Attention Is All You Need"~\citep{vaswani2017attention}. This document is followed by a proliferation of others conveying the same underlying knowledge: coding tutorials, blogs offering intuitive explanations, and forums clarifying ambiguities. Such texts are a \textit{natural byproduct of society's collective effort to communicate and process knowledge}.

We term these manifestations of the same underlying knowledge \textbf{auxiliary views}. We make this distinction because human-generated views are rarely paraphrases: they discuss the knowledge in diverse contexts (e.g., Transformers' advantages over other architectures) and forms (e.g., blogs, forums), collectively providing a more \textit{complete, contextualized} picture of $K$. We suspect pre-training is replete with these views, explaining its effectiveness beyond vague axes like quality or diversity. In our study, we frame paraphrasing as linguistic variation of a single view and set it as a control.

Our second research question concerns the knowledge represented around $K$. To acquire $K$, one may need prior knowledge it does not yet possess. We examine two types: \textbf{contextual knowledge} which $K$ references from the broader field, and \textbf{prerequisite knowledge}, the foundational concepts $K$ presupposes.


\subsection{Dataset} 

\textbf{Documents.}
We collect 36 documents across three domains to serve as $K$: twelve Computer Science papers from arXiv, twelve legal opinions from U.S. federal appellate courts, and twelve medical case reports from PubMed Central.  To mitigate confounding factors from OLMo2's pre-training corpus, we select documents published near or after the cutoff date and verify their absence through the Infini-gram API~\citep{liu2024infini}.  Additional details on sourcing, preprocessing, and the corpus-membership check are provided in Appendix~\ref{Appendix:preprocess}.

\textbf{Views.} For each document, we synthetically generate auxiliary views via LLMs; inspired by our earlier observation and prior work~\citep{gunasekar2023textbooks, allen2024physics, jiang2024instruction}, we construct textbooks, Stack Exchange--style Q\&A, and blogs (Table \ref{tab1}). Prerequisite knowledge is generated in the form of textbooks. Contextual knowledge is collected as cited arXiv papers and cited judicial opinions; medical case reports lack an analogous citation structure, so we omit contextual knowledge for this. Paraphrases are generated with GPT-4.1, while auxiliary and prerequisite views are generated with GPT-5 and GPT-5-mini. More details on data collection, pre-processing, and prompts for synthetic data generation are provided in Appendix~\ref{Appendix:preprocess}.

\begin{figure*}[t]
    \centering
    \includegraphics[width=\textwidth]{EMNLP/figures/probe_scaling_by_model_v13_mixed_reeval_v3_canonical.pdf}
    \caption{Knowledge acquisition across model sizes under a fixed token budget. Auxiliary views (green) improve both factual recall and inference over \textit{Source} (blue) and \textit{Para. 9} (orange); this advantage grows with model size.}
    \label{views}
\end{figure*}

\textbf{Probes}. To measure learning, we adopt LAMA-style probes~\citep{petroni-etal-2019-language,jiang-etal-2020-know, zhong2021factual}, where the task is to predict a sentence's final word or, following \citet{chang2024large}, a multi-word target. We design two probe types (Table~\ref{tab2}): factual probes, which measure recall of information stated explicitly in the text, and inference probes, which require combining one or more facts to infer information not stated. 

To construct factual probes, we filter for knowledge-bearing sentences, generate QA pairs from each, and convert them into self-contained, cloze statements; we automate this with GPT-5.4, having verified that it adequately understands the documents. The pipeline yields 6,435 factual and 430 inference probes, along with multiple-choice variants (4,515 factual and 322 inference MCQs), exceeding the scale of prior work~\citep{chang2024large}. We manually validate 100 probes of each type. Full pipeline details are in Appendix~\ref{Appendix:Probes}, and the dataset and code will be released publicly.\footnote{\url{pending}}

\subsection{Training Setup}

We use the base OLMo-2 models (1B, 7B, 13B, and 32B)~\citep{olmo20242} for training. Following \citet{chang2024large}, we employ single-batch knowledge injection: the documents fit in one forward pass, and the rest of the batch is filled with general data. We perform $N=100$ injections in our main experiments. In \textit{Source}, the original document is injected every batch. In \textit{Para.\ M}, we cycle through the original document and its $M$ paraphrases, repeating $\frac{N}{M+1}$ times. In \textit{Para.\ M + Aux.}, auxiliary views are injected alongside the paraphrased documents, with a blog, textbook chapter, and Stack Exchange Q\&A inserted into every batch. 

To remove the number of knowledge-bearing tokens as a confound, we token-match all conditions: since \textit{Para.\ M + Aux.} introduces additional knowledge-bearing tokens, we upsample the chunks in \textit{Source} and \textit{Para.\ M} so that every condition sees the same \# of tokens pertaining to $K$. Consequently, \textit{Para. M + Aux.} devotes fewer of its tokens to direct repetitions of the main document. Hyperparameters are in Appendix~\ref{append:hyperp}.


\textbf{General Data.} For general data replay, we stream tokens from the DCLM subset~\citep{li2024datacomp} of OLMo2's pre-training corpus.

\textbf{Metrics.} We evaluate the model's performance on our probes throughout training:
\begin{itemize}[itemsep=0pt, topsep=0pt, parsep=0pt]
    \item \textbf{Log-Prob.}: For cloze probes, we measure the joint log-probability of the target span.
    \item \textbf{Target Rank}: We measure the vocabulary rank of the correct target token under teacher forcing. For multi-token targets, we report the worst rank across the span; lower is better.
    \item \textbf{Multiple-Choice Accuracy}: For MCQ probes, we use 5-shot prompting with general-knowledge examples and constrain decoding to the answer choices.
\end{itemize}

\section{Large Language Models Learn Better with Auxiliary Views}
\label{aux}
\subsection{The Benefit of Auxiliary Views}

\textbf{Better Knowledge Acquisition.} In Figure~\ref{fig1:probe_32b}, even under a matched token budget, allocating tokens to auxiliary views substantially improves performance on both factual and inference probes. We see that this pattern holds across all of our metrics, including log probability, MCQA, and target rank (Figure~\ref{views_target_rank}). Note, we ensure all factors are held equal across conditions; the only difference is the injected auxiliary views replace a portion of direct views/paraphrases of the source document. We find that this ordering, \textit{Para.\ 9 + Aux} $>$ \textit{Para.\ 9} $>$ \textit{Source}, is established early during training and holds through convergence, though \textit{Source} learns faster on factual probes for the first $\sim$20 steps. 

Out of factual recall and inference, the improvement on inference is the more intuitive of the two. Although source documents alone should be sufficient to answer these probes too, the source offer only a single perspective. As Table~\ref{tab1} illustrates, these auxiliary views reformulate the knowledge in different styles and contexts, covering more ground on what the knowledge exactly entails. Learning with these views should lead to a more generalized encoding of knowledge, and thus inference probes are a natural place for the benefits to emerge.

\textbf{Generalization to Factual Recall.} Unexpectedly, auxiliary views also improve \textit{factual} recall, even though every factual target span is a verbatim phrase from the source. One would expect diverting tokens away from the source to \textit{hinder} direct memorization, not help it. This further suggests that auxiliary views encourage the model to encode the knowledge in a more generalized manner that, in turn, supports more effective factual recall \citep{spiro2017remembering}; we return to this in Section~\ref{sec:mechanistic} by looking at how this emerges mechanistically. 

\textbf{Model Size Effects.} The benefit of auxiliary views further \textit{emerges with scale}. As shown in Figure~\ref{views}, the 1B model derives little advantage from them, while the gap over the \textit{Source} and \textit{Para. 9} conditions widens steadily through 7B, 13B, and 32B. This trend is largely consistent across all domains (Figure~\ref{views_by_domain}). The lone exception is factual probes in the medical domain. We do not have a definitive explanation, especially as factual recall via MCQA returns to the overall trend. One possibility is that medical cases are more atomic (patient values, dosages) for which auxiliary views offers less benefits than the more conceptual arXiv and legal documents, and the probes are more sensitive to verbatim phrasing for such facts.

\textbf{Measuring Lexical Bias.} A possible confound is that synthetically generated auxiliary views leak the probes, especially for inference probes where the target span, unlike factual probes, need not appear in the source. We address the question of \textit{lexical} bias, whereby the model might favor particular phrases as the auxiliary views and probes were produced by related model families (GPT-5 and GPT-5.4, respectively). We measure how often each probe's target span occurs in the source, paraphrases, and auxiliary views.

Table~\ref{tab:target-occurrence-all} reports both frequency and coverage. Across all metrics, auxiliary views contain the target spans \textit{less} frequently and with \textit{lower} coverage than either the source or paraphrases. This strengthens our findings: the auxiliary views perform better on the probes despite stating the targets less.
% Requires \usepackage{booktabs,multirow} in the preamble.
\begin{table}[h]
\centering
\small
\setlength{\tabcolsep}{4pt}
\renewcommand{\arraystretch}{0.9}
\caption{Frequency and coverage of our probe targets across documents. \textit{Freq.} = occurrences per 1k tokens; \textit{Cover.} \ = fraction of targets present at least once across the documents. \textit{Para. 49} averages \textit{Freq} over 49 paraphrased documents but coverage includes presence in any of the paraphrases.}
\label{tab:target-occurrence-all}
\begin{tabular}{llrrrr}
\toprule
Probes & Corpus & \multicolumn{2}{c}{Full target} & \multicolumn{2}{c}{Bigram} \\
\cmidrule(lr){3-4}\cmidrule(lr){5-6}
 & & Freq. & Cover. & Freq. & Cover. \\
\midrule
\multirow{3}{*}{Factual} & Source & 0.391 & 0.70 & 0.829 & 0.92 \\
 & Para. 49 & 0.222 & 0.60 & 0.681 & 0.91 \\
 & Aux. & 0.112 & 0.33 & 0.437 & 0.74 \\
\cmidrule(lr){1-6}
\multirow{3}{*}{Inference} & Source & 1.225 & 0.51 & 0.703 & 0.71 \\
 & Para. 49 & 1.111 & 0.56 & 0.605 & 0.78 \\
 & Aux. & 1.070 & 0.51 & 0.628 & 0.77 \\
\bottomrule
\end{tabular}
\end{table}

\textbf{Human Auxiliary Views.} To test whether our findings depend on synthetic text, we repeat the experiment with human-written auxiliary views collected from the open web. The benefit again grows with model size (Figure~\ref{fig:human}), mirroring the synthetic-view results. This experiment covers only two documents, so we treat it as suggestive rather than conclusive; nonetheless, it provides  initial evidence that the effect is not merely an artifact of clean, synthetic text.



\begin{figure}[h]
    \centering
    \includegraphics[width=0.9\columnwidth]{EMNLP/figures/human.pdf}
    \caption{With human-written auxiliary views collected from the open web, the benefit grows with model size, consistent with our synthetic results. \(\Delta\) final log prob.\ is the gap between training with human auxiliary views and without (\textit{Para. 9}). Averaged over two documents.}
    \label{fig:human}
\end{figure}


\textbf{Which Auxiliary View?} In this experiment, we again make sure to token-match each of the conditions. Comparing each family of auxiliary views (textbooks, blogs, stackexchange) reveals that each of the views are similarly effective, but there is a synergistic effect where mixing multiple views performs slightly better, suggesting that diverse auxiliary views contribute to a more complete understanding (Figure \ref{views_ablation}). 



\subsection{Mechanistic Signatures of Auxiliary Views}
\label{sec:mechanistic}


\begin{figure}[h]
    \centering
    \includegraphics[width=\columnwidth]{EMNLP/figures/aux_base_minus_source_base_cosine_distance_all_projections.pdf}
    \caption{Per-channel difference in cosine distance from the base model, \textit{Para. 9 + Aux.} minus \textit{Source}, across layers and FFN channels (gate, up, down projections). Red indicates channels that auxiliary views move more than source; blue, less. A negative band around layers 16--24 shows that auxiliary views change the upper-middle layers \textit{less} than source, while affecting the middle and final layers \textit{more} .}
    \label{fig:mech_heatmap}
\end{figure}

Having established that auxiliary views improve knowledge acquisition, we ask how this phenomena emerges mechanistically. We compare each trained model's feed-forward network (FFN) weights against the base model along two axes: the \textit{magnitude} of change (relative delta norm and cosine distance) and its \textit{concentration} across channels (Gini coefficient), measured per layer and over training (Figures~\ref{fig:mech_layer},~\ref{fig:mech_steps}).

\textbf{Why MLP Channels?} We focus our analysis on the feed-forward (MLP) layers for two reasons. First, they account for the majority of a transformer's parameters, making them the natural locus for studying where knowledge is written during training. Second, they have comparatively clean interpretations: \citet{geva2021transformer} show that individual channels  in FFN layers operate as key--value memories. This channel-level view lets us read parameter change as movement in the model's memory, rather than as an opaque aggregate.

\textbf{Learning More by Changing Less.} Based on Figure~\ref{fig:mech_steps}, paraphrasing induces the \textit{largest} parameter movement of the three conditions, in both norm and cosine distance, yet adding auxiliary views \textit{reduces} this magnitude, though still more than source-only training. We interpret this as evidence that the two conditions learn differently. Paraphrasing supplies many surface-level variants of a single view, and the model appears to expend parameter change absorbing this lexical variation. Auxiliary views instead supply the knowledge in genuinely distinct framings, which the model can apparently integrate with less weight movement. This is consistent with our conjecture that auxiliary views enable a more generalizable encoding of knowledge: forming a more general, reusable representation requires overwriting fewer parameters than memorizing surface forms. Magnitude of change is not the same as quality of learning.

\textbf{Layer-wise Biases.} Based on Figure~\ref{fig:mech_layer}, parameter change is also not uniform across layers. All conditions concentrate change in the middle and final layers, with a pronounced dip in the upper-middle layers (Figure~\ref{fig:mech_layer}). Auxiliary views accentuate this structure: relative to source-only training, they change the middle and final layers \textit{more}, but in the upper-middle band (layers \(\sim\)16--24) they change parameters \textit{less than source}. Figure~\ref{fig:mech_heatmap} makes this bias visible: the delta in cosine distance (\textit{Para.9 + Aux.} \(-\) \textit{Source})  is negative in a band around layers 16--24 and positive elsewhere. Auxiliary views do not move more weight everywhere; they redistribute \textit{where} learning occurs, sparing a band of upper-middle layers while encoding more into the middle and final layers. 

 At 1B, where auxiliary views confer little benefit (Figure~\ref{views}), the corresponding per-channel difference is almost uniformly negative or none across all layers and projections (Figure~\ref{fig:mech_heatmap_1b}), lacking the layer-wise biases seen at 7B. The mechanistic signature of auxiliary views also emerges only at scale, mirroring the behavioral effect.

\begin{figure*}[t]
    \centering
    \includegraphics[width=\textwidth]{EMNLP/figures/batch_scaling_log_prob.pdf}
    \caption{\textbf{(Left)} Training OLMo2-7B with a batch size of 64 on \textit{Source} versus \textit{Para.~9}, tracking performance on inference probes on log probabilities. \textbf{(Right)} Final log probabilities of probes across model sizes and batch sizes regarding \textit{Source} versus \textit{Para.~9}. This experiment is on 6 documents to examine dynamics at smaller batch sizes.}
    \label{fig:batch_size}
\end{figure*}

\section{When Does Paraphrasing Help?}
\label{paraphrasing}

Paraphrasing is known to aid knowledge acquisition~\citep{allen2024physics, ovadia2024fine}, but prior reports conflict on its effectiveness. We revisit it to show that its benefit is conditional, and in doing so reconcile these reports.

\textbf{Paraphrasing prevents collapse.} Repeated exposure improves performance only up to $\sim$20 exposures, after which \textit{Source} saturates and then sharply degrades as the model overfits (Figure~\ref{fig:batch_size}, left). \textit{Para.\ 9} prevents this collapse, sustaining improvement through $\sim$40 exposures. As with auxiliary views, this benefit emerges only at 7B and above; at 1B, paraphrasing is slightly harmful.


\textbf{The benefit depends on batch size.} However, we see that this advantage depends on batch size, likely because larger batches mix in more general data per step and this similarly suppresses overfitting. As batch size grows, \textit{Source} becomes more stable and nearly matches \textit{Para.\ 9} by batch size 256, while \textit{Para.\ 9} stays roughly flat (Figure~\ref{fig:batch_size}, rightmost). Paraphrasing's inference gains thus come from preventing a degeneration that large batches also prevent; together, the two interventions become redundant.

Factual acquisition behaves differently. At small batch sizes, paraphrasing improves factual learning at 7B and above, whereas \textit{Source} sees almost no gain from larger batches (aside from a small improvement at batch size 64 for 7B). We attribute this advantage to how paraphrasing semantically varies the facts: although factual probes are drawn explicitly from source sentences, they are recast as self-contained atomic statements, so \textit{Source}'s verbatim memorization transfers less well. The advantage shrinks at larger batches and reverses at batch size 256 (7B and 13B), likely because the increased general-data mixing dilutes the gradient signal from the paraphrased text.


\textbf{Reconciling Prior Works.} This batch-size dependence reconciles conflicting prior work: \citet{chang2024large}, at batch size 2048 with 2048-token chunks, reports \textit{degraded} factual learning from paraphrasing, whereas \citet{allen2024physics}, at batch size 96 with 512-token chunks, reports \textit{gains}. The two sit at opposite ends of the regime we characterize here.

\section{Prior Knowledge Matters}

\begin{table}[h]
\centering
\begin{tabular}{@{} l c c @{}}
\toprule
\textbf{Model} & \textbf{Base} & \textbf{CPT} \\
\midrule
OLMo-2-0425-1B & 0.4380 & 0.5454 \\
OLMo-2-1124-7B & 0.6859 & 0.7272 \\
\bottomrule
\end{tabular}
\caption{The prior-knowledge gap. MCQA accuracy on prerequisite topics for each document, before (\textit{Base}) and after (\textit{CPT}) continued pre-training on synthetically generated textbooks covering that foundational knowledge.}
\label{tab4}
\end{table}

\textbf{Prior-knowledge Gap.} We first establish that the model lacks some of the foundational knowledge $K$ presupposes. We generate MCQA pairs on prerequisite topics for each document and measure the base model's accuracy. As shown in Table~\ref{tab2}, initial accuracy is low and rises after continued pre-training on synthetic textbooks covering this material, confirming a gap that domain adaptation must contend with. Recent work has begun to study how such gaps shape learning~\citep{wang2025objectivefinetuningllmsprior, gekhman2024does, yang2024fine}, but the question remains open.

\textbf{Distinct Benefits.} Adding either knowledge type on top of \textit{Para.\ 9} improves probe performance, but the two help different probe types (Figure~\ref{fig:prior_knowledge}). Contextual knowledge yields the larger gain on \textit{factual} probes in both domains, whereas prerequisite knowledge yields the larger gain on \textit{inference}. All conditions are token-matched as before.

To ask whether these gains reflect surface overlap, we measure how often each probe's target span appears in the inserted texts (Table~\ref{tab:target-occurrence-insert}), as in our earlier lexical-bias analysis. The factual gain is potentially a surface effect in part: cited works contain the factual targets about twice as often as the prerequisite textbooks (by coverage), so contextual knowledge's recall advantage plausibly reflects this greater lexical presence. The inference gain resists this explanation. Prerequisite knowledge does \textit{not} contain the inference targets more than contextual knowledge does---coverage is comparable and bigram frequency is in fact lower---yet it produces a substantially larger inference gain. This advantage therefore cannot be attributed to surface presence, and is more consistent with prerequisite knowledge supplying the foundations that let the model better integrate and infer new facts.

\section{Ablations}

\textbf{Learning Rate.} Across our experiments, we use a fixed peak learning rate of 4e-5. Following \citet{parmar2024reuse}, however, continued-pretraining recipes should inherit the base model's pre-training learning rate, and the 13B and 32B models were pre-trained at higher rates (9e-5 and 6e-5) than the 7B (3e-5)~\cite{olmo20242}. Because learning rate strongly governs how much is learned, holding it fixed at 4e-5 may have caused us to underestimate the model-size effect, and the effect of auxiliary views or paraphrasing overall. Consistent with this, ad-hoc analysis in Figure~\ref{fig:lr} shows that the advantage of auxiliary views and paraphrasing over source-only  widens as the learning rate increases.
 
\textbf{Order of Prior Knowledge.} In our experiments involving training on prerequisite concepts, one might suspect that the order in which the prerequisite concepts appear during training matters, akin to curriculum learning. We test this by varying whether prerequisite data is placed at the beginning, middle, or end of training. Counterintuitively, front-loading prerequisites performs slightly worse than placing them in the middle or end (Figure~\ref{fig:order}). That said, the differences are small. We leave a fuller investigation of curriculum effects in this setting to future work.

\section{Discussion \& Conclusion}

% We argue that the prevalence of auxiliary views in pre-training corpora has been a key factor in enabling LLMs to reach their present level of knowledge. In one way, this is unsurprising as humans suffer from the same problem in which we need help with each other across various domains and it's rather the human as a collective that has a full mastery not any individual.

% From a practical perspective regarding domain adaptation, our findings underscore the importance of careful data curation and, in some cases, targeted synthetic augmentation. For broad or foundational knowledge, pre-training corpora already contain abundant auxiliary views, particularly textbooks, which we find to be especially effective. In contrast, open scientific corpora rely heavily on primary research papers that lack such pedagogical scaffolding. Our results suggest that in these settings, augmenting training data with auxiliary views could be crucial for enabling models to acquire knowledge more effectively. Finding such data or synthetically generating these views for the latest scientific knowledge, however, may be a challenge.

% The findings of this research yield two primary contributions. First, we present what we believe to be the first to identify (a) the causal significance of auxiliary views, and (b) the distinct challenges associated with continued pre-training, notably including prior knowledge gaps, data replay, and post-instruction emergence. Addressing our core research question, we establish causal links between auxiliary view and enhanced knowledge encoding. These auxiliary views provide a robust explanation for the effectiveness of pre-training corpora beyond simple scale, suggesting that pedagogically structured texts are a critical determinant of successful knowledge acquisition and improved compositional reasoning. Further, our framework offers a robust pipeline for generating probes to extend this line of inquiry into other domains. Second, we offer practical guidelines for practitioners engaged in domain adaptation. Our recommendations include the strategic continuation of pre-training for foundational knowledge, in-batch strategy of data replay, and employing text paraphrasing techniques when data resources are scarce. In conclusion, our evidence underscores the crucial role of auxiliary views in fostering a more generalizable knowledge representation, leading to improvements in both factual recall and compositional reasoning. While acknowledging that the scope and scale of this study may limit full generalizability, we are confident that our hypothesis and supporting results will provoke deeper discussion and stimulate subsequent experimentation into the fundamental mechanisms of knowledge acquisition and intelligence within LLMs. Ultimately, we hope these insights prove valuable to practitioners in the domain adaptation space, encouraging a greater consideration for the importance of data.

\textbf{Auxiliary Views.} To our knowledge we are the first to isolate the significance of auxiliary views in knowledge acquisition during pre-training. The premise is not unsurprising: data augmentation has always been helpful, and diverse data has been a known principle in pre-training recipes. However, 



The human collective learns together, communicating to and teaching one another across varied contexts. Auxiliary views are simply the textual residue of this collective effort, and a model trained on them inherits its benefits.

in which they improved not only inference but, counterintuitively, factual recall. We show that this effect is not merely behavioral: auxiliary views leave a distinct mechanistic signature, moving fewer parameters and redistributing learning toward the middle and final layers. Together, these results offer an explanation for the effectiveness of pre-training corpora beyond simple scale, and lend a clearer mechanism to the commonly preached notion of "good" and "diverse" data.

\textbf{Practical Takeaways.} Second, we offer guidance for practitioners engaged in domain adaptation: continue pre-training on prerequisite knowledge to close foundational gaps; apply paraphrasing when data is scarce, while recognizing that its benefit is conditional on batch size and diminishes as batches grow; and consider targeted synthetic augmentation with auxiliary views. We find this last point especially relevant for scientific domains. Open scientific corpora are constantly evolving, producing primary research papers that lack the auxiliary views surrounding more established knowledge. Our results suggest that in these settings, synthesizing such views may be crucial for acquiring knowledge effectively; though doing so for the newest knowledge, which large language models may lack the domain expertise to generate high-quality auxiliary views for, remains a challenge.

Overall, we argue that the prevalence of auxiliary views in pre-training corpora has been a key factor in enabling LLMs to reach their present level of knowledge. In one sense this is unsurprising: humans face the same constraint. No individual holds full mastery of a domain; rather, it is the collective that does, through learning from and teaching one another across varied contexts. Auxiliary views are simply the textual residue of this collective effort, and a model trained on them inherits its benefits.

While the scope of this study may limit its generalization, we hope these insights prove valuable to practitioners in domain adaptation, encouraging greater consideration for how knowledge is represented, and provoke deeper discussion into the mechanisms of knowledge acquisition in LLMs.

\section*{Limitations}

Our analysis is restricted to three domains, which may introduce corpus-level biases. While our main results regarding auxiliary views were consistent across each domain, this is not the case for our secondary results regarding contextual and prerequisite knowledge. We acknowledge that this limits the generalization of our findings.

Our setup also introduces a batch scale limitation: our models are exposed to a limited amount of new knowledge at once, with batch sizes as high as 256, whereas pre-training involves simultaneously many more, with batch sizes usually no smaller than 2048 and usually spread very diverse domains instead of our controlled setup of injecting. While we believe our insights should scale under approximate independence between domains, it is possible that emergent behaviors arise at larger batch sizes or when learning more knowledge simultaneously. For instance, large-scale learning might give rise to unexpected multi-task synergies. Studying these effects at scale remains an important open question beyond the scope of this study. 

Furthermore, while our results carry important implications for understanding the nature of knowledge in LLMs, they are constrained by the scale of models we study (up to 32B), which is repeatedly emphasized in the literature as the primary determinant of capabilities in LLMs.

\bibliography{main}

\appendix
\section{Data}
\label{Appendix:preprocess}

\subsection{Dataset Construction}

\paragraph{Document collection.} The arXiv set was manually collected.  For medical documents, we use the NCBI E-utilities API to search PubMed Central for open-access case reports and download JATS XML.  For legal documents, we use the CourtListener REST API to search U.S. federal appellate opinions. Across domains, we filter candidate documents to obtain twelve target documents of suitable length, e.g., court opinions longer than 4000 tokens were rejected to avoid awkward chunking.

\paragraph{Pre-training corpus check.} We check all 36 source documents against the Infini-gram API.  For each document, we query the title and ten randomly sampled body sentences against the OLMo2 pre-training mixture index, \texttt{v4\_olmo-mix-1124\_llama} and found zero counts.

\paragraph{Cleaning.}
For arXiv papers, we clean the raw LaTeX by removing comments, figures, tables, presentation-only commands, unresolved includes, page breaks, and bibliography/appendix material; expanding simple macros; extracting title, abstract, and main body; resolving revision commands; and repairing erroneous line breaks while preserving LaTeX environments.  Medical case reports are converted from JATS XML into sectioned plain text.  Legal opinions require additional PDF/OCR cleanup: we remove page headers, docket/page artifacts, decorative separators, extracted footnotes, and merge dangling sentences into paragraphs.

\paragraph{Contextual Knowledge.}
In our study, contextual views are cited works associated with the source document.  For arXiv documents, we resolve each paper to an arXiv identifier, retrieve reference metadata using Semantic Scholar and OpenAlex, identify references with arXiv identifiers, download their source packages, and clean them with the same arXiv cleaning procedure.  For legal documents, we use CourtListener citation links to collect cited judicial opinions and apply the same legal-opinion cleaning procedure.  We do not construct contextual views for the medical documents because the case reports do not have cited works.

\subsection{Synthetic Data Generation}

\paragraph{Data statistics.}
Table~\ref{tab:dataset_artifact_inventory} summarizes the artifacts used in our experiments.  We generate 49 paraphrases per source document, yielding 1764 paraphrases in total.  For auxiliary views, $N$ counts the generated units obtained by splitting each view family into per-post blogs, per-question Stack Exchange Q\&A, and per-chapter textbooks. 



\begin{table}[h]
\centering
\scriptsize
\setlength{\tabcolsep}{3pt}
\begin{tabular}{@{}lrrr@{}}
\toprule
\textbf{Material} & \textbf{$N$} & \textbf{Avg. Len. (tokens)} & \textbf{Total Tokens} \\
\midrule
\multicolumn{4}{@{}l}{\textit{Overall}} \\
Papers & 36 & 5858 & 210{,}888 \\
Paraphrases & 1764 & 6180 & 10{,}904{,}160 \\
Blogs & 225 & 1901 & 427{,}632 \\
Stack Exchange Q\&A & 453 & 1290 & 584{,}392 \\
Textbook Chapters & 282 & 2239 & 631{,}403 \\
Prerequisite Chapters & 582 & 4282 & 2{,}492{,}392 \\
Contextual Documents & 639 & 10416 & 6{,}655{,}888 \\
\midrule
\multicolumn{4}{@{}l}{\textit{Computer Science}} \\
Papers & 12 & 10550 & 126{,}598 \\
Paraphrases & 638 & 11002 & 7{,}019{,}384 \\
Blogs & 105 & 1920 & 201{,}648 \\
Stack Exchange Q\&A & 242 & 1460 & 353{,}321 \\
Textbook Chapters & 158 & 2872 & 453{,}833 \\
Prerequisite Chapters & 187 & 4555 & 851{,}724 \\
Contextual Documents & 506 & 10791 & 5{,}460{,}148 \\
\midrule
\multicolumn{4}{@{}l}{\textit{Medical}} \\
Papers & 12 & 3765 & 45{,}181 \\
Paraphrases & 588 & 3858 & 2{,}268{,}311 \\
Blogs & 72 & 1979 & 142{,}500 \\
Stack Exchange Q\&A & 117 & 1068 & 124{,}973 \\
Textbook Chapters & 61 & 1404 & 85{,}666 \\
Prerequisite Chapters & 220 & 4589 & 1{,}009{,}582 \\
Contextual Documents & 0 & -- & -- \\
\midrule
\multicolumn{4}{@{}l}{\textit{Legal}} \\
Papers & 12 & 3259 & 39{,}109 \\
Paraphrases & 589 & 3274 & 1{,}928{,}141 \\
Blogs & 48 & 1739 & 83{,}484 \\
Stack Exchange Q\&A & 94 & 1129 & 106{,}098 \\
Textbook Chapters & 63 & 1459 & 91{,}904 \\
Prerequisite Chapters & 175 & 3606 & 631{,}086 \\
Contextual Documents & 133 & 8991 & 1{,}195{,}740 \\
\bottomrule
\end{tabular}
\caption{Dataset statistics by text type and domain.}
\label{tab:dataset_artifact_inventory}
\end{table}

\paragraph{Paraphrases.}
Paraphrases are generated from cleaned source documents on a paragraph-level basis.  We preserve section headers verbatim and avoid paraphrasing LaTeX-only paragraphs. Domain-specific prompts are used for academic, legal, and medical text.  The academic prompt preserves LaTeX formatting, equations, proper nouns, titles, section headers, and technical terminology; the legal prompt preserves legal meaning, party names, citations, quoted language, dates, docket numbers, and legal terms of art; and the medical prompt preserves diagnoses, medications, doses, routes, timelines, lab values, imaging findings, procedures, outcomes, and clinical terminology. We use GPT-4.1 to generate paraphrases with a temperature of 1 and a top-p of 0.975.

\paragraph{Auxiliary views.}
For each document, we generate synthetic auxiliary views conditioned on the target document and intended to restate, explain, or pedagogically reorganize the same document-level content.  The Computer Science setting generates textbook chapters for college students with basic machine-learning background, Stack Exchange--style questions from a confused student followed by grounded answers, and technical blog posts for a broader technical audience.  The medical setting adapts these formats to clinical education: case-based textbook sections for medical students or residents, clinical teaching Q\&A, and clinical blog posts.  The legal setting uses casebook/treatise-style chapters for law students, Law Stack Exchange-style Q\&A, and analytic legal commentary.  For each view family, we first generate an outline, list of questions, or list of blog ideas, and then generate the individual chapters, answers, or posts from that plan. We use GPT-5 to generate the outlines and GPT-5-mini to generate the contents of the chapters, questions, or posts.

\paragraph{Prerequisite views.}
Prerequisite views are generated separately from auxiliary views.  For arXiv papers, we use a curriculum-design prompt to create prerequisite textbook chapters that teach the foundations needed to understand the paper while excluding the paper's own novel ideas.  For medical case reports, the prompt asks for general medical background--for example relevant anatomy, pathophysiology, pharmacology, diagnostic interpretation, differential diagnosis, and standard management--while explicitly forbidding patient-specific chronology, workup, treatment course, outcome, or novel observations from the source case.  For legal opinions, the prompt similarly asks for doctrinal and procedural background while excluding the source case's parties, facts, outcome, holding, and novel reasoning. When useful for legal background, we additionally include landmark or doctrinally foundational cited opinions as context for the generated prerequisite chapters. We use GPT-5 to generate the outlines and GPT-5-mini to generate the contents of the chapters.



\FloatBarrier

\section{Probe Generation}
\FloatBarrier
Our probe construction pipeline is adjusted and refined for each domain (arXiv, legal, and medical). We provide the pipeline for arXiv documents. An overview is shown in the Figure~\ref{figA1} (next page). The full pipeline for all domains will be made available in the to be released code.
\begin{figure*}[t]
    \centering
    \includegraphics[width=0.95\textwidth]{overview-SFig1.pdf}
\caption{Overview of the Probe Generation Pipeline for Arxiv Documents. \textbf{(A) Factual Probe Generation}: We generate probes on a sentence-by-sentence basis to adhere closely to the source document for ground truth. \underline{Preprocessing}: We first filter sentences that contain minimal knowledge (to prevent noise from structural comments, e.g., "Let us first discuss the following results"). We then use heuristics to filter out sentences that may yield low-quality probes, such as those that are too short (e.g., "the sweep has 22 runs") or contain excessive \LaTeX{} code with no valid English targets for extraction. \underline{Question Extraction}: From each filtered sentence, we extract 1--3 questions that capture the knowledge contained within. \underline{Contextualization}: Questions are contextualized to be clear and self-contained, accurately reflecting the knowledge being tested. \underline{Cloze Conversion}: The questions are converted into cloze statements, placing the answers at the end. \underline{Refinement}: We refine the statements, ensuring all mathematical content is written in \LaTeX{} and verifying the correctness of the preceding steps. \textbf{(B) Compositional Probe Generation}: Recognizing that not all atomic facts necessitate a compositional probe (some only concern simple experimental details), we employ a two-level approach. First, we break the paper into sections and prompt a Large Language Model (LLM) to extract compositional questions. Compositionality is defined as either (1) \textbf{Inference} (using reasoning upon supporting text to reach a new insight) or (2) \textbf{Synthesis} (combining knowledge across several facts into a synthesized statement). This is performed section-by-section to ensure granularity. Second, we repeat this step by providing the entire paper to the LLM to allow for more holistic questions. \underline{Cloze Conversion}: Similar to the factual process, these questions are converted into cloze statements with the answer at the end. \underline{Refinement}: Statements are refined to ensure mathematical content is in \LaTeX{} and to enforce proper self-containment by referencing the paper (e.g., "according to the paper"). \underline{Filtering}: A comprehensive filtering step is conducted, where we ask the LLM to assess if the questions are too simple, imprecise, or confusing. This step, on average, removes 15\% of probes.}
    \label{figA1}
\end{figure*}

% \FloatBarrier


% \section{Source-Weighted BM25+}\label{Appendix:BM25+}
% To quantify the lexical advantage present in the synthetic views (e.g., paraphrases, textbooks) we generate over the \textit{Source}, we propose a \textbf{Source-Weighted BM25+} metric. Traditional IDF weighting is unsuitable here, as our goal is not to rank documents but to measure the lexical advantage over \textit{Source}. Instead, we substitute IDF with a weight $w_{\text{source}}(t)$ derived from the \textbf{Source} term frequency ($\text{tf}_{\text{source}}$). This prioritizes terms that are rare or absent in the original papers but frequent in the synthetic views, effectively capturing the lexical ``leakage'' of our probes.

% For a probe $P$ consisting of $n$-gram terms $t \in P$, and an auxiliary training corpus $C$, the transfer relevance score is defined as:

% \begin{equation}
%     \text{Score}(P, C) = \sum_{t \in P} w_{\text{source}}(t) \cdot \left( \frac{\text{tf}_{C}(t) \cdot (k_1 + 1)}{\text{tf}_{C}(t) + k_1 \cdot \left(1 - b + b \cdot \frac{|D_C|}{\text{avgdl}_C}\right)} + \delta \right)
% \end{equation}

% where $\text{tf}_{C}(t)$ is the term frequency of $t$ in the auxiliary corpus $C$, and $k_1$, $b$, and $\delta$ are standard BM25+ hyperparameters (set to $k_1=10.0$, $b=0.75$, $\delta=1.0$). 

% \begin{equation}
%     w_{\text{source}}(t) = \frac{5}{1 + \text{tf}_{\text{source}}(t)}
% \end{equation}

% This weighting scheme ensures that terms which are common in the Source contribute negligibly to the score, while terms that are rare in the Source are weighted highly.

% \FloatBarrier

% \section{Data Replay}
% % \FloatBarrier

% \begin{figure*}[t]
%     \centering
%     \includegraphics[width=0.75\textwidth]{data_replay_1B_7B_comparison_exposure_steps.pdf}
% \caption{Comparison of factual-probe performance when training OLMo2-7B on \textit{Source} versus \textit{Para.~9} with a batch size of 32. Paraphrasing yields modest improvements when contrasted with the larger gains observed on compositional probes.}
%     \label{fig:data_replay}
% \end{figure*}
% \FloatBarrier
% \section{Hits@K}
% \FloatBarrier

% \begin{figure*}[t]
%     \centering
%     \includegraphics[width=\textwidth]{hits_source.pdf}
%     \caption{Hits@k progression during training on \textit{Source} across model sizes. For clarity, we restrict the figure to target spans that appear in \textit{Source}, removing any probes absent from the training distribution to avoid confounding effects from inherently difficult or unseen items. Despite extensive training, it is difficult to do well on these probes.}
%     \label{fig:hits_k_source}
% \end{figure*}

% \begin{figure*}[t]
%     \centering
%     \includegraphics[width=\textwidth]{hits_para.pdf}
%     \caption{Hits@k progression during training on \textit{Para. 9} across model sizes with a batch size of 64. For clarity, we restrict the figure to target spans that appear in \textit{Para. 9}, removing any probes absent from the training distribution to avoid confounding effects from inherently difficult or unseen items. Despite extensive training, it is difficult to do well on these probes.}
%     \label{fig:hits_k_para}
% \end{figure*}
% \FloatBarrier

% \section{Additional Plots for Paraphrasing and Repeating Data}
% \FloatBarrier

% \begin{figure*}[t]
%     \centering
%     \includegraphics[width=0.75\textwidth]{factual_scaling.pdf}
% \caption{Comparison of factual-probe performance when training OLMo2-7B on \textit{Source} versus \textit{Para.~9} with a batch size of 32. Paraphrasing yields modest improvements on factual probes when contrasted with the larger gains observed on compositional probes.}
%     \label{fig:factual_source_para}
% \end{figure*}

% \begin{figure*}[t]
%     \centering
%     \includegraphics[width=0.75\textwidth]{batch_scaling_hits.pdf}
%     \caption{Training on \textit{Source} for 500 exposures with cosine warmup and no decay. We observe an early dip followed by steady increases in target-span log-probabilities. Hits@10 improves as well, though the gains are comparatively modest and experience no initial dip.}
%     \label{fig:batch_scaling_hits}
% \end{figure*}


\FloatBarrier
\clearpage
\twocolumn[{
\section{Additional Plots}
% 
% \begin{figure*}[h]
%     \centering
\begin{center}
    \includegraphics[width=\textwidth]{EMNLP/figures/probe_scaling_by_domain_v13_mixed_reeval_v3_canonical.pdf}
    \captionof{figure}{Under a fixed token budget, auxiliary views lead to better generalization on both factual recall and inference; we see that the effect size of auxiliary views also increase with model size. We see that this pattern holds even when averaging by domain, though the factual probes for the medical documents are the sole exception. }
    \label{views_by_domain}
% \end{figure*}

% \begin{figure*}[h]
%     \centering
    \vspace{1em}

    \includegraphics[width=\textwidth]{EMNLP/figures/explanation_source_comparison_7b_reeval_v5_delta_all_domains.pdf}
    \captionof{figure}{Under a fixed token budget, each family of auxiliary views (textbooks, stackexchange QA, blogs) performs similarly; we see that mixing these views does slightly better on both factual and inference probes, implying synergistic benefits though minimal.}
    \label{views_ablation}
% \end{figure*}
\end{center}
}]
\clearpage
 
\begin{figure*}[!t]
    \centering
    \includegraphics[width=0.9\textwidth]{EMNLP/figures/probe_scaling_by_model_target_rank.pdf}
    \caption{Under a fixed token budget, auxiliary views lead to better generalization on both factual recall and inference; we see that the effect size of auxiliary views also increase with model size. Lower ranks is better.}
    \label{views_target_rank}
\end{figure*}

\begin{figure*}[!t]
    \centering
    \includegraphics[width=\textwidth]{EMNLP/figures/7b_parameter_delta_mlp_time_comparison.pdf}
\caption{FFN parameter change over training (OLMo2-7B). Paraphrasing induces the largest parameter movement, while auxiliary views reduce this magnitude toward source-only training. Both the concentration of change (Gini) and downstream performance converge by step \(\sim\)50 (performance not shown), whereas magnitude (relative delta norm, cosine distance) continues to grow afterward; parameters keep drifting after learning has effectively converged.}
    \label{fig:mech_steps}
\end{figure*}

\begin{figure*}[!t]
    \centering
    \includegraphics[width=\textwidth]{EMNLP/figures/7b_parameter_delta_mlp_layer_comparison.pdf}
\caption{Per-layer FFN parameter change from the base model over training (OLMo2-7B). \textit{Magnitude} (relative delta norm, cosine distance) and \textit{concentration} (Gini) are shown for the gate, up, and down projections. All conditions concentrate change in the middle and final layers with a dip in the upper-middle layers; auxiliary views increase change in the middle and final layers but reduce it relative to source in the \(\sim\)16--24 band.}
    \label{fig:mech_layer}
\end{figure*}

\begin{figure*}[!t]
    \centering
    \includegraphics[width=0.85\textwidth]{EMNLP/figures/1B_delta_parameters.pdf}
\caption{The same per-channel difference in cosine distance from the base model (\textit{Para.~9 + Aux.}\ minus \textit{Source}) as Figure~\ref{fig:mech_heatmap}, but for the 1B model. Unlike the structured pattern at 7B, the difference is either blue (negative) or white, indicating that auxiliary views change parameters slightly \textit{less} than source nearly everywhere, with no organized layer-wise band. The only exceptions are a few isolated channels in \texttt{down\_proj} at the deepest layers. }
\label{fig:mech_heatmap_1b}
\end{figure*}

\begin{figure*}[!t]
    \centering
    \includegraphics[width=\textwidth]{EMNLP/figures/7b_lr_sweep_para49_main_metrics.pdf}
\caption{Effect of peak learning rate (2e-5, 4e-5, 8e-5) on knowledge acquisition for the 7B model. We report final log-probability on factual and inference probes (left two panels) and final accuracy on factual and inference MCQA (right two panels), comparing source documents alone (\textit{Source}), 49 paraphrases (\textit{Para.\ 49}), and 49 paraphrases plus auxiliary views (\textit{Para.\ 49 + Aux.}). Across all metrics, the advantage of auxiliary views and paraphrases grows with learning rate.}
\label{fig:lr}
\end{figure*}

\begin{figure*}[!t]
    \centering
    \includegraphics[width=0.8\textwidth]{EMNLP/figures/prior_knowledge_match_histogram_vs_para_7B.pdf}
\caption{Gain in \(\Delta\) log prob.\ over \textit{Para.\ 9} from adding prerequisite versus contextual knowledge, by domain and probe type. Medical is omitted (no citation structure). Contextual knowledge helps factual probes more; prerequisite knowledge helps inference probes more.}
\label{fig:prior_knowledge}
\end{figure*}

\begin{figure*}[!t]
    \centering
    \includegraphics[width=0.9\textwidth]{EMNLP/figures/prior_knowledge_position_logprob_histogram_delta_all_domains_7B.pdf}
\caption{Effect of prerequisite-concept ordering on knowledge acquisition. We compare placing prerequisite data at the \textit{Front}, \textit{Middle}, or \textit{End} of training, reporting the change in log-probability ($\Delta$ Log Prob) on factual (left) and inference (right) probes. Ordering has only a minor effect, with \textit{Front} placement performing slightly worse than \textit{Middle} or \textit{End}.}
\label{fig:order}
\end{figure*}

 
\FloatBarrier
\section{Additional Tables}

\begin{table}[h]
  \centering
  \small
  \setlength{\tabcolsep}{4pt}
  \renewcommand{\arraystretch}{0.9}
  \caption{Frequency and coverage of our probe targets across the contextual and pre-requisite knowledge. \textit{Freq.} = occurrences per 1k (words for the full target, OLMo
  tokens for bigrams); \textit{Cover.}\ = fraction of targets present at least once (full
  target) or the mean fraction of a target's bigrams present (bigram). \textit{Prerequisites} = generated textbook chapters; \textit{Cited Works} = cited papers or legal opinions.}
  \label{tab:target-occurrence-insert}
  \begin{tabular}{llrrrr}
  \toprule
  Probes & Corpus & \multicolumn{2}{c}{Full target} & \multicolumn{2}{c}{Bigram} \\
  \cmidrule(lr){3-4}\cmidrule(lr){5-6}
   & & Freq. & Cover. & Freq. & Cover. \\
  \midrule
  \multirow{2}{*}{Factual} & Prerequisites  & 0.031 & 0.12 & 0.129 & 0.54 \\
   & Cited Works & 0.039 & 0.24 & 0.141 & 0.72 \\
  \cmidrule(lr){1-6}
  \multirow{2}{*}{Inference} & Prerequisites  & 0.827 & 0.21 & 0.154 & 0.47 \\
   & Cited Works & 0.838 & 0.25 & 0.095 & 0.54 \\
  \bottomrule
  \end{tabular}
  \end{table}
  
\FloatBarrier
\section{Hyperparameters for Replication and Details for Reproduction}\label{append:hyperp}
For training, we use the TRL library. Unless otherwise noted, our experiments use the following defaults: 
learning rate $4\times10^{-5}$, context size 4096, batch size 256, weight decay 0.1, 
cosine decay scheduler with 0.1 warm-up ratio and minimum learning rate ratio of 0.1, seed 42, max gradient norm of 1, and the AdamW optimizer 
($\beta_1 = 0.9$, $\beta_2 = 0.999$, $\epsilon = 10^{-8}$) with BF16 training. 


\section{Prompts for Probe Construction (Arxiv)}\label{Appendix:Probes}

We provide some of the prompts used in the our pipeline to construct our factual and inference probes for the arXiv documents. The complete set of prompts, which we also adjust for each domain (legal, medical), can be found in our code.%\syang{knowledge probe, inference probe? please revise the text here and the prompts descriptions} 

\textbf{1. Prompt to extract atomic facts from a paragraph:}
\lstinputlisting[style=prompt]{prompts/atomic_facts.txt}

\textbf{2. Prompt to extract atomic facts from a paragraph with context:}
\lstinputlisting[style=prompt]{prompts/atomic_facts_contextual.txt}

\textbf{3. Prompt to generate inference questions from the text:}
\lstinputlisting[style=prompt]{prompts/comprehension_questions.txt}


\section{Prompts for Synthetic Data Generation (Arxiv)}\label{Appendix:Prompts} %(paragraph, prior knowledge, explanations)}
% We paraphrase each arxiv paper by asking the LLM with temperature=1 and top\_p=0.95.

\textbf{1. Prompt to synthesize Stack Exchange style question-answer pairs:}
\\Question generation:
\lstinputlisting[style=prompt]{prompts/stackExchange1.txt}

Answer generation:
\lstinputlisting[style=prompt]{prompts/stackExchange2.txt}

Latex formatting refinement:
\lstinputlisting[style=prompt]{prompts/stackExchange3.txt}

\textbf{2. Prompt to synthesize textbook style explanations:}
\\Textbook outline generation:
\lstinputlisting[style=prompt]{prompts/textbook1.txt}

Chapter generation:
\lstinputlisting[style=prompt]{prompts/textbook2.txt}

\textbf{3. Prompt to synthesize blog post style explanations:}
\\Blog post ideas generation:
\lstinputlisting[style=prompt]{prompts/blog1.txt}

Blog posts generation:
\lstinputlisting[style=prompt]{prompts/blog2.txt}

\textbf{4. Prompt to generate prior knowledge chapters:}
\\List of prerequisite chapters generation:
\lstinputlisting[style=prompt]{prompts/prior1.txt}

Chapter generation:
\lstinputlisting[style=prompt]{prompts/prior2.txt}

\end{document}
